\section{Acquiring an exact smaller formulation}
\label{active:section}

The preceding formulation bounds concern a representation once it is known.
Removing redundant constraints can require information about the input.
We give families in which this acquisition cost, the resulting barrier,
and its metric can all be calculated. The search bounds are classical
quantum-query results; the point is their precise connection with an exact
conic reformulation. In particular, an existential small barrier need not
provide a cheap way to evaluate that barrier from the supplied data.

\subsection{Output and access conventions}

A source query returns one indexed input symbol, either classically or
coherently in the usual reversible oracle. Public arithmetic and the
preparation of index registers do not count as source queries. Their gate
and bit costs are separate. All bounded-error statements require success
probability at least $2/3$ on the complete output, on every promised input.
An \emph{exact reusable compiler} produces a finite classical description
of the exact projected body, usable without further source queries.
It can, for example, return its retained inequalities and a barrier for
their intersection. A wrapper that continues calling the source oracle
does not meet this output convention. Lower bounds also cover a classical
description with exact membership or boundary evaluation, provided that
this evaluation requires no further source queries. Computation using the
description may be arbitrarily expensive in a query lower bound.

\begin{lemma}[Acquisition of independent marked positions]
\label{active:search}
Let $b\geq1$ and $m\geq2$. In each of $b$ public blocks of $m$ bits,
exactly one bit is one. Returning all marked positions has quantum query
complexity $\Theta(b\sqrt m)$ and randomized query complexity
$\Theta(bm)$. If exactly $k$ of $N$ bits are marked, with $k$ known,
all marks can be found with certainty in $O(\sqrt{Nk})$ quantum queries.
\end{lemma}
\begin{proof}
Exact amplitude amplification with a known number of marked positions
finds a mark in $O(\sqrt{n/t})$ queries when $t>0$ of $n$ positions are
marked \cite[Theorem~4]{BHMT2002}. Apply this separately to the blocks.
For known total cardinality, suppress previously found indices and sum
$O(\sqrt{N/t})$ for $t=k,k-1,\ldots,1$. There is no accumulation of
failure probabilities in these exact-query upper bounds.

For completeness, index block inputs by $r\in[m]^b$ and use the
nonnegative adversary matrix
\[
 \Gamma=\sum_{j=1}^b
 I^{\otimes(j-1)}\otimes(J_m-I_m)\otimes I^{\otimes(b-j)}.
\]
Its norm is $b(m-1)$. A query to position $\ell$ of block $j$ retains
only the $j$th summand, replacing $J_m-I_m$ by the adjacency matrix
of the star joining $\ell$ to the other positions. Its norm is
$\sqrt{m-1}$. The spectral adversary bound therefore gives
$\Omega(b\sqrt{m-1})=\Omega(b\sqrt m)$; this is the standard
direct-sum argument \cite[Theorems~3--4]{ACGT2010}.
Classically, under independent uniform marked positions, success
probability $2/3$ for all outputs implies that probability for each
individual position. Revealing the other blocks in advance cannot
make a uniform marked position in an unread block less uncertain.
A successful search in that block requires $\Omega(m)$ expected
queries, even allowing a final unqueried guess. Summing over blocks
gives $\Omega(bm)$; scanning gives the matching upper bound.
\end{proof}

The exact quantum upper bounds use the standard query model, which permits
the public rotations needed for exact amplitude amplification. An
approximate gate implementation must separately account for its precision.

\subsection{Nested constraints and consensus}

Let $d\geq1$ be fixed and let $r_i=2-b_i$, where
$b\in\{0,1\}^N$ has weight zero or one. Consider
\begin{equation}\label{active:consensus}
 \max e_1^Tz_1,\qquad
 (r_i,z_i)\in Q_{d+1},\quad z_i=z_{i+1}\quad(1\leq i<N).
\end{equation}
Its projection is exactly $\{z:\norm z\leq r_*\}$, where
$r_*=\min_i r_i$, and its optimum is $r_*$. The origin is strictly
feasible. The public path equalities have constant incidence; only the
right-hand-side radii are hidden.

\begin{proposition}[Exact radius compilation]\label{active:radius}
The exact projected ball and an optimizer of
\eqref{active:consensus} can be acquired with $\Theta(\sqrt N)$
quantum queries or $\Theta(N)$ randomized queries. The same lower
bounds hold for an optimum-value estimate with additive error below
$1/3$. After acquisition,
\[
 f_{r_*}(z)=-\log(r_*^2-\norm z^2)
\]
is a barrier with exact parameter one. On the all-radius-two input,
the natural restricted product barrier is $Nf_2$ and has exact
parameter $N$. Its central path from $z=e_1$ to
$z=(2-\epsilon)e_1$, $0<\epsilon\leq1/2$, has length
$\Theta(\sqrt N\log(1/\epsilon))$, compared with
$\Theta(\log(1/\epsilon))$ for $f_2$.
\end{proposition}
\begin{proof}
The acquisition problem is the zero-versus-one search promise.
Exact amplitude amplification, followed by verification of the returned
index, decides it in $O(\sqrt N)$ queries
\cite[Theorem~16]{BHMT2002}. The search lower bound is
\cite{BBBV1997}; the uniform single-mark distribution gives the
classical lower bound. The optimum values are separated by one.
The fixed ball barrier has parameter one by the ball calculation
in Section~\ref{sec:concrete-barriers}. Alternatively, direct inversion of
its Hessian gives gradient squared dual norm
$2\norm z^2/(r_*^2+\norm z^2)<1$, with boundary limit one;
self-concordance follows from its Lorentz restriction.
Multiplication by the integer $N$ preserves self-concordance and
multiplies this exact parameter by $N$. On the axis,
\[
 f_2''(a)=\frac{2(4+a^2)}{(4-a^2)^2},
\]
whose square root is comparable to $(2-a)^{-1}$ on $[1,2)$.
The length follows by integration. Reflection symmetry and strict
convexity put every center for objective $e_1$ on this axis.
\end{proof}

Every move of local norm at most $R\in(0,1)$ has metric length at most
$-\log(1-R)$. Hence the displayed natural-barrier central schedules
require $\Omega_R(\sqrt N\log(1/\epsilon))$ moves. This is a statement
about that barrier and those local schedules. One cannot obtain a
uniform replacement simply by dividing the original barrier by $N$.
If exactly one radius is one and $N>1$, its normal restriction near
the tighter boundary is
$-N^{-1}\log(1-a)+O(1)$ with smooth bounded corrections.
For $-c\log t$, the self-concordance ratio
$|f'''|/[2(f'')^{3/2}]$ equals $c^{-1/2}$. The divided barrier
therefore violates self-concordance near that boundary.

There is a separate conditioning issue in the uneliminated consensus
system. At its analytic center and in one coordinate, put
$D=\operatorname{Diag}(2/r_i^2)$ and let $A$ be the path incidence
matrix. Then
\begin{equation}\label{active:path-kkt}
 K=\begin{pmatrix}D&A^T\\A&0\end{pmatrix},\qquad
 \kappa(AD^{-1}A^T)=\Theta(N^2),\qquad
 \kappa(K)=\Theta(N^2).
\end{equation}
Here condition numbers use singular values and $N\geq2$.
Indeed, $I/2\preceq D\preceq2I$, whereas the extreme eigenvalues of
$AA^T$ have orders $N^{-2}$ and one. For a negative eigenvalue
$-t$ of $K$, eliminating its primal component gives
$A(D+tI)^{-1}A^Ty=t y$. The smallest positive solution has order
$N^{-2}$ by these same comparisons and the min--max principle;
all positive eigenvalues are bounded below by $1/2$, and all
eigenvalues have bounded magnitude. Thus $K$ has the stated
condition number despite its forest support. Public consensus
elimination removes this factor. Equation~\eqref{active:path-kkt}
is consequently not a lower bound on the cost of a linear-system
algorithm allowed that elimination.

More generally, suppose $N$ constraints are partitioned into $k$
supplied chains under inclusion, of lengths $n_1,\ldots,n_k$.
An ordered queried key identifies each chain's smallest member.
Once that member is known, suppose it has a specified barrier of
parameter $\nu_j$. Minimum finding in each chain
\cite{DurrHoyer1996}, amplified to joint success, gives
\begin{equation}\label{active:chains}
 O\!\left(\log(2k)\sum_{j=1}^k\sqrt{n_j}\right)
 \leq O\!\left(\sqrt{Nk}\log(2k)\right)
\end{equation}
queries and a compiled barrier with parameter
$\nu_0+\sum_j\nu_j$ when an additional public barrier of parameter
$\nu_0$ is present. Common strict feasibility and positive definiteness
on the effective affine space are understood. This bound covers
homothetic copies $r_i C$ of a public compact convex body containing
zero: nestedness leaves only the minimum radius.

The dependence $\sqrt{Nk}$ is necessary in general. For $N=km$,
take $k$ independent interval chains $|z_j|\leq2-b_{j\ell}$,
with each chain containing zero or one mark. Returning all optimum
coordinates of $\max\sum_jz_j$ within $1/3$ returns all $k$ OR
values. The direct sum of the $k$ star adversaries has norm
$k\sqrt m$, while a query filters it to norm one. The lower bound
is $\Omega(k\sqrt m)$ and the classical bound is $\Omega(km)$.
For this promised family the upper bound is also
$O(k\sqrt m)$, by the exact promised search in the preceding proof;
no logarithmic amplification is needed. On the all-zero input the
natural and compiled barriers are respectively
$m\sum_j f_2(z_j)$ and $\sum_j f_2(z_j)$.
Their common diagonal path from zero to total objective gap
$\epsilon$, for $0<\epsilon\leq k/2$, has lengths
\[
 \Theta\!\left(\sqrt N\log(k/\epsilon)\right),\qquad
 \Theta\!\left(\sqrt k\log(k/\epsilon)\right).
\]
Additional public coupling may make the post-compilation optimization
nonseparable. Its solution cost is additional to
\eqref{active:chains}; the acquisition bound alone does not estimate it.

\subsection{A two-dimensional family with incomparable active factors}

For $v>0$ write
\[
 q(v)=(1+v,3-v^2),\qquad
 E(q)=\{z\in\R^2:q_1z_1^2+q_2z_2^2\leq1\},\qquad
 u_i=\frac{i}{2(N+1)}.
\]
The parameters used below lie in $(0,1]$, so both coefficients
are positive. A public sentinel is $E(q(1))=E(2,2)$.
An input trit at position $i$ selects $q(1)$, $q(u_i)$, or
$q_i^-=(1+u_i,1)$, denoted respectively by $0,+,-$.
Let $S$ be the plus positions. Since the sentinel dominates all
minus factors, the projected body is
\begin{equation}\label{active:pareto-body}
 D_S=E(q(1))\cap\bigcap_{i\in S}E(q(u_i)).
\end{equation}
Each factor is a transformed three-dimensional Lorentz slice.
A path of copies of the two coordinates gives constant incidence
and a scalar KKT graph of bounded treewidth: a bag can contain the
two coordinates of consecutive copies, their equality multipliers,
and one constant-size local cone block. Consecutive bags share
only a copy. The origin is strictly feasible.

\begin{proposition}[Pareto compilation]\label{active:pareto}
If $1\leq k\leq N/2$ and $|S|=k$ is promised, the $k+1$
factors of \eqref{active:pareto-body} are essential in that intersection
description. Exact reusable compilation has quantum query complexity
$\Theta(\sqrt{Nk})$ and randomized query complexity $\Theta(N)$.
These lower bounds concern an exact body description, not solution of
one fixed objective or the minimum number of factors in every possible lift.
\end{proposition}
\begin{proof}
For each retained $v$, put $s_v=(2v,1)$. The identity
\[
 \ip{q(v)-q(w)}{s_v}=(v-w)^2
\]
shows that squared coordinates
$s_v/\ip{q(v)}{s_v}$ define a point on factor $v$ and strictly
inside all the others. Thus removal of $v$ enlarges the body:
a small outward perturbation still satisfies the other inequalities.
It also follows that different sets of retained parameters give
different bodies. An exact source-free membership description can
distinguish them: evaluate it at the finitely many public witnesses
obtained by these perturbations for the universe of all parameters.
The perturbations may be chosen so that each witness violates just
its own factor in that universe.

The upper bound is the known-cardinality search of
Lemma~\ref{active:search}. For the lower bound restrict to one plus
position in each of $k$ public blocks of size
$m=\lfloor N/k\rfloor\geq2$, with all other positions zero.
The reusable description determines every marked position, so the
lemma applies. Since $km\geq N/2$, its bounds have the claimed orders.
\end{proof}

The retained constraints can be simultaneously active. If $v<w$
are consecutive retained parameters and $s=(v+w,1)$, then
\[
 \ip{q(r)-q(v)}s=(r-v)(w-r).
\]
Thus squared coordinates $s/\ip{q(v)}s$ give exactly the $v$ and
$w$ factors tight. Their outward normals are linearly independent;
a positive combination exposes this point. Its open normal cone
therefore supplies objectives using both factors. The projected
problem is a coupled planar intersection, not independent scalar
minimum problems.

Here the cost of a barrier can change much more than the dimension.
Set $F_q(z)=-\log(1-z^T\operatorname{Diag}(q)z)$.
On the hard branch with only zeros and pluses the natural and compiled
barriers are
\begin{equation}\label{active:pareto-barriers}
 F_o=M F_{q(1)}+\sum_{i\in S}F_{q(u_i)},\qquad
 F_c=F_{q(1)}+\sum_{i\in S}F_{q(u_i)},\qquad M=N-k+1.
\end{equation}
Their exact parameters satisfy
$M\leq\nu(F_o)\leq N+1$ and $\nu(F_c)\leq k+1$.
For the lower bound approach the sentinel boundary on the positive
first axis. The other factors and all their derivatives stay bounded,
and the squared gradient dual norm tends to $M$. The upper bounds
follow by adding the parameter-one ellipsoid barriers.
The planar body itself admits a barrier of parameter two
\cite{LeeYue2021}; the $M$ lower bound is specific to $F_o$.

For objective $\max z_1$, both central paths have $z=(a,0)$,
$0\leq a<2^{-1/2}$. If $f_A(a)=-\log(1-Aa^2)$, then
\[
 f_A''(a)=\frac{2A(1+Aa^2)}{(1-Aa^2)^2}.
\]
For marked factors, $1<A<3/2$, so $f_A''$ is bounded above
and below by positive universal constants on this whole interval.
The sentinel square-root Hessian is comparable to
$(2^{-1/2}-a)^{-1}$ near its endpoint. Integrating the square
root of the Hessians in \eqref{active:pareto-barriers}, and using
$\max(\sqrt x,\sqrt y)\leq\sqrt{x+y}\leq\sqrt x+\sqrt y$,
gives, uniformly for $0<\epsilon\leq1/4$,
\begin{equation}\label{active:pareto-lengths}
 L_o(\epsilon)=\Theta\!\left(\sqrt M\log(1/\epsilon)\right),\qquad
 L_c(\epsilon)=\Theta\!\left(\sqrt k+\log(1/\epsilon)\right).
\end{equation}
Lengths run from the origin to $a=2^{-1/2}-\epsilon$.
The assumption $k\leq N/2$ ensures $M\geq k$.
The endpoint estimate of Lemma~\ref{lem:cm-height} also gives
an $\Omega(\sqrt M\log(1/\epsilon))$ distance lower bound for
$F_o$ at sufficiently small universal $\epsilon$: its barrier
increase is at least $M\log(c/\epsilon)$, while
$\nu(F_o)\leq M+k\leq2M$. Thus bounded local steps incur the
corresponding lower bound even if the path leaves the central curve.

\subsection{A growing-dimensional fixed-objective decoder}

The distinction between body compilation and a fixed objective can be
removed by allowing the projected dimension to grow. Let $N=km$,
$m\geq2$, and $u_\ell=\ell/[2(m+1)]$. In block $j$ exactly one
position $\ell_j$ is marked. In coordinates $(y,p_j,q_j)$ retain
the public sentinel and the marked inequality
\begin{equation}\label{active:growing-body}
 4(y^2+p_j^2+q_j^2)\leq1,\qquad
 y^2+(1+u_{\ell_j})^2p_j^2+(3-u_{\ell_j})^2q_j^2\leq1.
\end{equation}
Before compilation there are $m$ hidden factors in each block,
each either the sentinel or the marked factor, and one additional
public sentinel. The coefficient matrices are rational with
$O(\log m)$ bits, and all factors are direct $Q_4$ slices.
The common $y$ is the sole coupling between blocks.

\begin{proposition}[Fixed-objective acquisition and metric separation]
\label{active:growing}
The body \eqref{active:growing-body} has dimension $2k+1$ and
$2k$ essential retained inequalities. It has a lift with $\Theta(N)$
constant-size cone factors, bounded row and column incidence, and
bounded scalar KKT treewidth, all with bounds independent of $k,m$.
Both exact reusable compilation and a coordinate-readable solution
of $\max\sum_jq_j$ with coordinate error below $1/[40(m+1)]$
require and suffice with $\Theta(\sqrt{Nk})$ quantum queries or
$\Theta(N)$ randomized queries. The same lower bounds hold for
an exactly feasible coordinate-readable solution with total objective
gap below $1/[40(m+1)]$.

For the separate objective $\max y$, the natural and compiled
explicit barriers have parameters of orders $N$ and $k$ and central
lengths, from their common center to $y=1/2-\epsilon$,
\[
 \Theta\!\left(\sqrt N\log(1/\epsilon)\right),\qquad
 \Theta\!\left(\sqrt k\log(1/\epsilon)\right)
 \quad(0<\epsilon\leq1/4).
\]
\end{proposition}
\begin{proof}
The origin is strict. To make the marked factor of block $j$ alone
tight, set $y=p_j=0$, $q_j=(3-u_{\ell_j})^{-1}$, and all other
coordinates zero. Its sentinel is strict since $3-u_{\ell_j}>5/2$.
To make its sentinel alone tight, take $y=q_j=0,p_j=1/2$;
the marked factor is strict since $1+u_{\ell_j}<3/2$.
These points also prove essentiality after small outward perturbations.
For bounded incidence, give every slot a triple copy and join triples
along paths within blocks. Connect representative $y$ copies between
blocks along a path. A constant-size bag for consecutive copies,
their multipliers, and each local cone block gives a tree
decomposition; at block junctions only a constant number of
representative coordinates need be retained. This also proves the
assertion for the scalar KKT graph, including dense local cone Hessians.

For the decoding objective, every feasible point has
$q_j\leq(3-u_{\ell_j})^{-1}$. All these bounds are attained
simultaneously only at
\[
 y=0,\quad p_j=0,\quad q_j^*=(3-u_{\ell_j})^{-1}.
\]
The optimum is unique. Neighboring candidate $q_j^*$ values differ
by more than $1/[18(m+1)]$. A coordinate error below
$1/[40(m+1)]$ is less than half that spacing, so rounding recovers
every mark. For a feasible point the deficits $q_j^*-q_j$ are
nonnegative; a total objective gap below the same threshold bounds
each deficit. Lemma~\ref{active:search} proves the lower bounds and
the acquisition upper bounds. Once the marks are known, the exact
optimizer takes $O(k)$ arithmetic operations. The conclusion is for
readable coordinates, not the objective value alone or a normalized
quantum state.

For $\max y$, reflection symmetry gives $p_j=q_j=0$ on both
central paths. Write $f_A(y)=-\log(1-Ay^2)$. Their restrictions are
exactly
\[
 F_o(y)=Nf_4(y)+kf_1(y),\qquad
 F_c(y)=kf_4(y)+kf_1(y).
\]
The term $f_1''$ is uniformly bounded on $[0,1/2)$, whereas
$\sqrt{f_4''}$ has order $(1/2-y)^{-1}$ near the endpoint.
Integration proves the length assertions. The squared gradient dual
norm on this axis tends to $N$ or $k$, respectively, while
additivity gives upper parameters $N+k$ and $2k$.
The diagonal symmetry removes mixed Hessian entries on the axis,
so these are also lower bounds for the full barriers.
\end{proof}

The post-compilation problem is only coupled through one scalar.
Its explicit decoding objective is easy once the hidden constraints
are known. This is an acquisition advantage; it does not require
a quantum linear-system solver.

\subsection{Blockwise relative-entropy compilation}

The same acquisition phenomenon occurs for nonsymmetric cones.
Use $D(U\|V)=U\log(U/V)$, with its closed perspective extension.
Let $N=Bm$, $B\geq1$, $m\geq2$, and choose public values
$\gamma_1<\cdots<\gamma_m$ in $[1,2]$ with minimum spacing
$h\asymp m^{-1}$. Each block has exactly one marked position
$r_j$. Set $d_{ji}=1$ except $d_{j r_j}=e^{-\gamma_{r_j}}$.
An oracle query returns the mark bit; the coefficient at a marked
position is a public function of that position. Coefficients may
be retained symbolically, so no exact transcendental-number
encoding is hidden in this query convention. All source terms in
block $j$ share $U_j,V_j$.

The elementary identity \eqref{eq:qre-compiler} becomes
\begin{equation}\label{active:entropy-identity}
 \sum_{i=1}^m D(U_j\|d_{ji}V_j)
 =mU_j\log(U_j/V_j)+\gamma_{r_j}U_j
 =D(mU_j\|m e^{-\gamma_{r_j}/m}V_j).
\end{equation}
One can equivalently retain the rational/public coefficient
$\gamma_{r_j}$ in the middle expression. Projection of $m$
separate epigraphs, under $t_j=\sum_i t_{ji}$, is exactly the
single epigraph of \eqref{active:entropy-identity}. One implication
follows by summing inequalities. Conversely, distribute the
nonnegative aggregate slack among the $m$ epigraph coordinates.
This argument includes the closed perspective boundary whenever
the aggregate value is finite; if it is infinite, neither epigraph
has a point over those $U_j,V_j$.

\begin{proposition}[Entropy acquisition and matched-gap lengths]
\label{active:entropy}
An exact reusable compiler for these $B$ aggregate epigraphs has
query complexities
\[
 Q=\Theta(B\sqrt m)=\Theta(\sqrt{NB}),\qquad
 R=\Theta(Bm)=\Theta(N).
\]
The same lower bounds apply to a coordinate-readable, exactly
feasible solution of the slice $U_j=V_j=1$ minimizing
$\sum_jt_j$, with objective excess below $h/3$.
On that slice, at matched positive total gap $\Delta$, the
natural and compiled projected central points coincide. Between
gaps $\Delta_0>\epsilon>0$ their exact barrier-metric lengths are
\begin{equation}\label{active:entropy-lengths}
 \sqrt N\log(\Delta_0/\epsilon),\qquad
 \sqrt B\log(\Delta_0/\epsilon),
\end{equation}
for their respective displayed product barriers.
\end{proposition}
\begin{proof}
The projected epigraph in block $j$, evaluated at $U_j=V_j=1$,
has lower boundary $t_j=\gamma_{r_j}$. Hence an exact reusable
description determines all $r_j$; exact membership at the public
midpoints between candidate boundary values suffices.
Lemma~\ref{active:search} gives both bounds and the upper algorithm.
For the solution task, feasibility gives
$e_j=t_j-\gamma_{r_j}\geq0$. Total excess below $h/3$
implies $e_j<h/3$ for every $j$, and nearest-value rounding
recovers all marks.

With $U_j=V_j=1$, the usual relative-entropy cone barrier restricts
to $-\log(t_{ji}-D(1\|d_{ji}))$, up to a constant.
At multiplier $\eta>0$, all $N$ source slacks equal $1/\eta$,
so their total gap is $N/\eta$ and each projected block slack is
$m/\eta$. The compiled barrier has $B$ slacks, each
$1/\widehat\eta$, total gap $B/\widehat\eta$.
At $\eta=N/\Delta$ and $\widehat\eta=B/\Delta$, both project
to $t_j=\gamma_{r_j}+\Delta/B$.
A positive orthant logarithmic barrier has line element
$\sum_i(ds_i/s_i)^2$. Along these curves it is respectively
$N(d\Delta/\Delta)^2$ and $B(d\Delta/\Delta)^2$, proving
\eqref{active:entropy-lengths}.
\end{proof}

For the natural formulation the first length is measured in its
full source-slack coordinates. It is also the length of the common
projected path for the marginal barrier obtained by minimizing
over the source slack allocations: that marginal is
$-m\sum_j\log(t_j-\gamma_{r_j})$, up to a constant.
Equal multipliers do not give equal projected points; matched gaps
are essential. The statement does not extend here to moving
$U_j,V_j$. Allowing feasibility violations also changes the decoder:
negative deficits must be charged in addition to the objective excess.

The two full ambient exponential-cone products have optimal barrier
parameters $3N$ and $3B$, even among coupled barriers, by
Theorem~\ref{thm:entropy-product}. On the fixed slice in
Proposition~\ref{active:entropy}, the displayed parameters are $N$
and $B$. Neither ambient number is an intrinsic lower bound for
every barrier on the projected epigraph. The small relative-entropy
barriers of \citet{FS2023} and the positive-map constructions of
\citet{HeSaundersonFawzi2026} are the relevant prior barrier results.
Our identity is their elementary scalar setting, not a new
relative-entropy barrier theorem. The conclusions here concern
acquiring the effective coefficients, decoding readable outputs,
and comparing the specified metrics at equal gap.

For scalar variables, bounded incidence and bounded KKT treewidth
can be retained before compilation. Arrange each block as a binary
tree; carry copies of $U_j,V_j$ down its edges and sum the
epigraph variables up the same tree. Each node has only a constant
number of variables and equations, and each parent--child separator
has constant size. These separators give an explicit bounded-width
tree decomposition. For matrix entropy variables the separator size
grows with the matrix dimension; no dimension-free scalarized
treewidth claim follows.

\subsection{What the oracle advantage includes}

Consensus reconstruction is public: a projected vector supplies every
copied coordinate. For a nonzero amplitude-encoded vector, equal copying
uses the public map $|z\rangle\mapsto|u_N\rangle|z\rangle$.
An explicit array of all $Nd$ copied entries still costs
$\Omega(Nd)$ output operations. In the entropy example an implicit
online compiler could instead search one requested block in
$O(\sqrt m)$ queries and retain access to the input; that output
does not incur all-block acquisition until all blocks are requested.

There is also no free explicit-input loading theorem. With bounded
fan-in $g>1$, a readable OR output sensitive to $m$ source bits
has a backward dependency cone containing those bits, and therefore
requires $\Omega(m)$ input incidences and depth
$\Omega(\log_g m)$. Applying this to the independent promised OR
chains gives total $\Omega(N)$ incidences and depth
$\Omega(\log_g(N/k))$. Parallel classical reductions match these
orders. The query improvements above consequently assume a supplied
coherent oracle or an already constructed access structure.

Redundancy removal before an interior-point solve is classical
presolve \cite{Gondzio1997}; quantum search on partially ordered
sets is also established \cite{Montanaro2009}. Our chain model gives
unordered access to inclusion keys, whereas the poset-search models
in the latter reference have their own comparison or ordering
promises. Minimum finding, amplitude amplification, and adversary
direct sums are established methods. The explicit families
above combine them with essential-constraint witnesses, sparse conic
lifts, readable-output reductions, and exact metric comparisons.
The results distinguish that conjunction from a general quantum
optimization speedup, an all-barriers lower bound, or a universal
claim that fewer retained inequalities always give the best lift.
